\chapter{Remaining Work and Timeline}
\label{ch:remaining}

Everything in this chapter is planned work; none of it is built or measured yet unless stated. The milestones follow
the project plan (\code{CLAUDE.md}), and are taken one at a time, each with a short plan approved before it starts.

\section{Open items from M2 and M3}

\begin{itemize}
  \item \emph{Fix A}: make the evidence field of \code{submit_answer} explicitly required for a ``yes'' answer, and
        measure whether the cycle-check rejections disappear (all 18 rejections of the M2 run were of this kind).
  \item \emph{Fix B}: state in the system prompt that $G$ is undirected (read from the graph), and measure the effect
        on the node-degree questions without the \code{degree} tool.
  \item \emph{Fix C}: fewer text tool calls (a clearer prompt or an example call), measured on strict accuracy.
  \item A run on the small graphs, skipped in the M2 closing run (only large graphs were run).
  \item Reading the \code{combine16k} results into the learning log, and rerunning the combine setups with a second
        model at 16k context.
  \item \emph{M3 step 6, close-out}: the last step of M3. Like every milestone, it ends with a short entry in the
        learning log.
  \item Further items from the GDS Agent review: batch tools that accept handles as inputs (for example the degrees of
        all nodes in a stored result, or the top-$k$ nodes by a metric among candidates), error-triggered one-line
        hints, node names as identifiers with ``did you mean'' suggestions, a ``do not retry'' note in truncation
        messages, and a count of how often frontier models actually verify their results in the published GDS traces.
\end{itemize}
\todo[inline]{Confirm the exact scope of M3 step 6 (close-out); the plan only names it.}

\section{Model pilot}

Before further experiments, the cheap model is to be fixed by a pilot of four models that fit the 12~GB GPU
(Table~\ref{tab:candidates}): qwen3.5:9b, gemma4:12b, gpt-oss:20b and qwen3:8b, each on 54 large-graph questions (six
per task, one run), about 5--10 minutes per model on the RTX 5070. The decision criteria are strict accuracy, rescued
tool calls, verifier rejections, and tokens and time per question; the newest small models in the Ollama library
should be checked first. The chosen model and tool set are then fixed for the remaining experiments.

\section{M4: router and skills}
\label{sec:m4}

The router classifies the question into a task family (which also fixes the answer type and the verifier) and
extracts the parameters (node identifiers). It is meant to be swappable, and four candidates are to be compared on
routing accuracy, latency and cost: a small LLM with structured output; TypeSafe Jev, a zero-shot classifier with typed
output and a confidence score (a candidate for the classification part only, as a cloud API and a new dependency);
embeddings (\code{nomic-embed-text}); and a regular-expression baseline. Low confidence leads to no verifier
(``unverified'') or a question to the user, never a guessed verifier. Jev is a router candidate only, never a verifier.

Skills are per-family playbooks, ported from the pseudo-code prompts of the prior work (\code{data/pseudocodes/}).
Following the GDS Agent experience, they should use progressive disclosure: a short skill, with troubleshooting text
shown only when an error happens. A skill ablation (none, minimal, prescriptive; per family and per model) is planned,
since the code hint helped one model a lot. With the router in place, the agent sees only its family's tools
(principle 6), which shortens every prompt.

\section{M4.5: interactive command-line interface}

An interactive interface on top of the router: load any graph file, ask questions in plain English in a loop, and see
the tool calls and whether the answer was verified. Later, follow-up questions that reuse earlier answers. It serves
debugging and the thesis demonstration.

\section{M5: escalation}

A model cascade: when the verifier rejects an answer (or the run ends without one), a stronger model reruns the loop.
The local candidate for the stronger model is qwen3.6:35b-a3b, a mixture-of-experts model with about 3B active
parameters that needs expert offload to system memory; hosted models are the alternative. Escalation can only be
triggered by verifiable failures, so its benefit depends on the share of tasks with evidence.

\section{M6: benchmarks and baselines}

\paragraph{Generator.} The question generator of the prior work, ported and scaled to graphs with 10,000 nodes and
more, with balanced yes/no answers, controlled density (long paths, not every pair an edge), questions for
topological sort, bipartiteness and the multi-step family (farthest node, $k$-hop counts), and later weights and
direction.

\paragraph{Benchmark adapters.} Loaders for GABench \cite{gabench}, ProGraph \cite{prograph}, GrAlgoBench
\cite{gralgobench} and GT Bench \cite{gtbench}. A London adapter for the GDS Agent benchmark \cite{gdsagentbench} is a
priority: the London Underground graph loads into networkx without Neo4j, and the benchmark already contains frontier
results with tokens and dollar cost per question, so the expensive baseline is already paid for. Some answers depend
on how GDS implements an algorithm (PageRank scaling, seeded community detection), so only questions with exactly one
correct answer are usable; the review estimates roughly 55--60 of the 88, to be checked question by question with
every reference answer recomputed. The questions name stations, so node names are needed, and models may inject
outside knowledge about London, which synthetic graphs avoid.

\paragraph{Baselines.} Pure prompting (the prior work's methods, with the graph serialized in the prompt), a generic
ReAct-style agent with Python, Claude Code headless, and GraphTeam \cite{graphteam}, under the rule that harnesses are
compared with the same model and, where possible, the same tools. The MCP server provides the tools; serving
\code{submit_answer} over MCP is needed so that a general harness returns a gradable, structured answer. Sandboxing for
the Claude Code and Codex baselines is open (a local microVM sandbox is a candidate).

\section{M7: experiments}

\begin{itemize}
  \item Size scaling, from the dataset sizes to 10,000+ nodes, where serialization stops being possible and handles
        matter.
  \item Accuracy-versus-cost Pareto fronts across models and harnesses (tokens, latency and dollars).
  \item Component ablations: tools (fitting tools versus low-level tools), verifiers, handles, code execution, router
        and skills, escalation, compaction.
  \item At least three runs per configuration, bootstrap confidence intervals, and a majority-answer baseline per task.
  \item Final numbers on the test split (graphs 50--99), which has not been used so far.
  \item Process-level evaluation (Option B of the proposal) as a possible chapter: process metrics, trajectory
        classes, and whether they predict answer accuracy.
\end{itemize}

\section{Timeline}

\todo[inline]{Add dates for each milestone (pilot, M3 close-out, M4, M4.5, M5, M6, M7, writing) and the submission deadline. The repository does not record a schedule.}

\begin{table}[ht]
  \centering
  \small
  \caption{Milestones and status on 9 October 2026.}
  \label{tab:milestones}
  % source: CLAUDE.md (Milestones, Current status)
  \begin{tabularx}{\textwidth}{l>{\raggedright\arraybackslash}X l}
    \toprule
    Milestone & Content & Status \\
    \midrule
    M0 & Project setup, LiteLLM smoke test & done \\
    M1 & Bare loop, five tools, \code{submit_answer}, JSONL logging & done \\
    M2 & Verifiers for the ten problems; accuracy and cost report & done (2026-10-07) \\
    M3 & Full tool set, MCP server, input validation, handles, sandbox & steps 1--5 done; close-out open \\
    pilot & Choice of the cheap model & planned \\
    M4 & Router and skills; per-family tool filtering & planned \\
    M4.5 & Interactive CLI & planned \\
    M5 & Escalation / model cascade & planned \\
    M6 & Benchmark adapters and baselines & planned \\
    M7 & Experiments: scaling, Pareto fronts, ablations, repeated runs & planned \\
    \bottomrule
  \end{tabularx}
\end{table}
